\section{Conclusion}
\label{sec:conclusion}

We have proposed a checker that makes editing a list while anything else still holds the old version a
compile error in beni, a pure language compiled to JavaScript, so that every list can be a plain array
written in place with no run-time ownership information and no annotations in user code. The property is
uniqueness at the update site, decided by liveness; the inference is a directional, field-sensitive cell
analysis with per-path liveness and four side conditions, whose function summaries are polymorphic over
the ownership classes of function-typed parameters and are solved for recursive groups by an optimistic,
validated fixpoint. The design depends on three properties of the language---strict evaluation in source
order, saturated calls with no automatic currying, and an existing inference of per-function facts
published in interfaces---and on ten assumptions about the runtime, the library and the platforms, the
most demanding of which is that a UI runtime hand the model over to \code{update} and keep nothing of the
old one.

We believe the soundness argument has the right shape and the right assumptions, and we have tried to say
where it is weakest: the entry state of \code{update}, the platform obligations, and the exclusive-contents
flag on which nested edits rely. What we cannot know from the design alone is whether the rule is worth
its cost. It rejects correct programs---every program that keeps an old version, every list captured by a
command that a later message edits---in exchange for a guarantee about performance, and its benefit to the
two applications we examined depends on widening the set of operations that write in place, which would
reject more. The next step is the report-mode experiment of \S\ref{sec:evaluation}: build the analysis,
run it over real programs, classify every rejection by hand, and decide on the numbers.
